\section{Related Work}\label{sec:related_works}

\subsection{Query Containment}\label{sec:related_works_qc}

Query containment is the problem of deciding, for every possible database, whether the result of a query $Q_1$ is a subset of the result of a query $Q_2$, denoted $Q_1 \sqsubseteq Q_2$; in this case $Q_1$ is the \emph{contained query} and $Q_2$ the \emph{containing query}.
Query containment has been used in multiple query optimization tasks, including view selection, caching, and query rewriting~\cite{Afrati2019-ch2, Shirkova2011, Halevy2001}.

Query containment under bag or bag-set semantics is a strictly stronger relationship than its set-semantic counterpart~\cite{Chaudhuri1993}.
The literature does not describe general necessary conditions for containment under bag-set and bag semantics~\cite{ioannidis1995containment, khamis2021bag}; however, sufficient conditions have been identified.
Furthermore, the general case of bag or bag-set semantic query containment for \ac{UCQ} has been shown to be undecidable~\cite{ioannidis1995containment, Jayram2006, Afrati2010}.
However, some fragments of bag-set and bag semantic conjunctive queries have been proven decidable~\cite{khamis2021bag, KOPPARTY20111097, Chaudhuri1993, Afrati2010}.
\Cref{tab_decidable_containment_bag_bag_set} summarizes the bag and bag-set \ac{CQ} containment fragments relevant to this paper and their complexity.\footnote{For the general case, containment mappings give a sufficient but not a necessary condition.}

\begin{table}
  \caption{Complexity of bag and bag-set \ac{CQ} containment for the fragments relevant to this paper.}
  \label{tab_decidable_containment_bag_bag_set}
  \begin{tabular}{lll}
    \toprule
    Case & Semantic & Complexity\\
    \midrule
    General \ac{CQ}~\cite{Chaudhuri1993} & Bag, bag-set & $\Pi^p_2$-hard\\
    \ac{CQ}, $Q_1$ proj.-free~\cite{Afrati2010} & Bag-set & NP-complete\\
    \ac{CQ}, $Q_1,Q_2$ proj.-free~\cite{Afrati2010} & Bag-set & $O(n^2\log n)$\\
    \ac{CQ}, $Q_1,Q_2$ proj.-free~\cite{Afrati2010} & Bag & $O(n^2\log n)$\\
    \bottomrule
  \end{tabular}
\end{table}

Beyond relational databases, query containment has been studied in other graph query languages, including XPath~\cite{schwentick2004xpath} and Cypher~\cite{tang2025proving}.
In the context of SPARQL, set-semantic query containment has received considerable attention~\cite{WudageChekol2013, letelier2013static, chekol2012sparql, chekol2012sparql2, jenaQC, Spasic2023}.
Extensions incorporating RDF schemas, OWL entailment regimes~\cite{WudageChekol2013} and property paths~\cite{Kostylev2025, chekol2016containment} have also been proposed, though these fall outside the scope of this paper.


\subsection{Federated Queries}\label{sec:related_works_fed}

Federated SPARQL queries have been approached from two perspectives: explicit federation through the \texttt{SERVICE} clause, and automatic source selection.
The SPARQL 1.1 Federated Query specification~\cite{w3SPARQLFederated} introduces the \texttt{SERVICE} clause as the W3C standard mechanism for directing subexpressions to specific endpoints.
The SPARQL specification mentions the use of variables as endpoint \acp{IRI} in \texttt{SERVICE} clauses but does not give an official interpretation of their behavior.\footnote{\url{https://www.w3.org/TR/sparql11-federated-query/\#variableService}}
\citeauthor{BuilAranda2013}~\cite{BuilAranda2013} introduced \emph{Service-Safeness}, a syntactic rule outside the SPARQL specification that makes such queries computable by guaranteeing a finite set of federation members.
Because Service-Safeness is not part of the specification and multiple SPARQL query engines use their own interpretation, this paper does not consider variables as endpoint \acp{IRI}.
\citeauthor{Cheng2021}~\cite{Cheng2021} proposed FedQPL, a formalism to express federated queries under automatic source selection.
Automatic source selection algorithms\rt{Let's cite them here! E.g. FedX, ... I'd just include all algorithms that were mentioned and cited in the paper of Elias. We can also include Comunica in this list.} alleviate the practical difficulty of determining which federation member each subquery should be sent to~\cite{hanski2025observations}.
Thus, users can query a federation without knowing how the data is distributed.
Furthermore, \citeauthor{Cheng2021}~\cite{Cheng2021} introduce \ac{ESA}, a canonical source-selection algorithm that requires no prior information about the federation.
In \ac{ESA}, each triple pattern is evaluated against every federation member and the results are combined via bag union, guaranteeing complete results regardless of how data is distributed across the federation.

To the best of our knowledge, no contribution has addressed query containment for federated SPARQL queries; this is the gap our paper addresses.
